\documentclass[11pt,a4paper]{article}
\usepackage[T1]{fontenc}
\usepackage[utf8]{inputenc}
\usepackage[margin=1in]{geometry}
\usepackage{amsmath,amssymb}
\usepackage{booktabs}
\usepackage{url}
\usepackage[hidelinks]{hyperref}
\hypersetup{pdftitle={What a Metamorphic Relation Can and Cannot See in an FHE Decoder},
            pdfauthor={Bader Alissaei},
            pdfsubject={Fully homomorphic encryption, metamorphic testing},
            pdfkeywords={CKKS, FHE, metamorphic testing, conformance, decoder}}
\usepackage{microtype}

\title{What a Metamorphic Relation Can and Cannot See in an FHE Decoder\\[2mm]
\large A degree criterion for scale errors, and why three laws missed one}

\author{Bader Alissaei\\ \small VaultBytes Innovations Ltd\\ \small \texttt{b@vaultbytes.com}}
\date{}

\begin{document}
\maketitle

\begin{abstract}
Metamorphic and relational testing checks an implementation against identities its outputs should
satisfy, which is attractive for fully homomorphic encryption because it needs no reference
implementation and no trusted golden vectors. We give a criterion for what such a test can see. A
relation assembled from device decodes is blind to a multiplicative decoder error exactly when it is
invariant under $D \mapsto sD$, and it is blind to an additive one exactly when its coefficients sum
to zero. The two conditions are independent, and a suite can satisfy both without anyone noticing,
because neither is visible in the identities as usually written. We apply the criterion to a CKKS
conformance suite of our own and find that all three of its laws are blind to a scale error, so a
decoder scaled by $1.01$ passes with a higher score than an honest one. The repair needs no oracle:
the multiplicative homomorphism mixes degrees in $D$ and its residual is $s(1-s)D(ab)$, which we
measure. We also show that sensitivity is not refusal. At the floor our own reference runs at, that
law scores $6.66$ bits against a floor of $6.0$ and still returns a pass.
\end{abstract}

\section{The question}

A test that needs no reference implementation is worth having. For fully homomorphic encryption it
is worth more than usual, because a reference implementation presumes both a trusted party and a
trusted channel, which is the third party an accelerator buyer is trying to avoid, and because a
check that decrypts user data defeats the point of the scheme. So relational testing is the natural
design, and recent work takes it: Eidolon fuzzes FHE libraries with noise-aware mutation over
equivalent arithmetic expressions~\cite{eidolon}, and HERTA applies metamorphic relations derived
from FHE semantics across the software stack and reports 21 previously unknown
bugs~\cite{herta}.

The question this note answers is narrow. Given a suite of such relations, which faults can it see?
We answer it for the decoder, which is where an FHE implementation turns ring elements back into
numbers and where an error is both easy to introduce and consequential in every answer the engine
will ever give.

\section{The criterion}

Write $D$ for the decoder a correct implementation realises and $D'$ for the one a faulty engine
realises. Take the faulty decoder to be affine,
\[
D'(x) = s\,D(x) + \delta,
\]
which covers a scale error, a constant offset, and any combination. Let the tested relation be
\[
\sum_i a_i\, D(x_i) = 0,
\]
an identity that holds for correct decoding, with the $x_i$ produced by the device and the $a_i$
fixed by the relation. Substituting the faulty decoder gives the residual
\[
s \sum_i a_i D(x_i) \;+\; \delta \sum_i a_i \;=\; \delta \sum_i a_i ,
\]
since the first sum vanishes by the identity. Two consequences follow, and they are independent.

\paragraph{The additive error.} The residual is $\delta \sum_i a_i$, so the relation sees $\delta$
exactly when the coefficients sum to something other than zero. This is often paraphrased as both
sides decoding the same number of times, and for relations whose coefficients are $\pm 1$ that
paraphrase is correct. It is not correct in general. The relation $D(c \boxdot a) = c\,D(a)$ for a
public scalar $c$ decodes once on each side, yet its coefficients sum to $1 - c$, so it does see
$\delta$ whenever $c \neq 1$. The criterion is the coefficient sum, not the decode count.

\paragraph{The multiplicative error.} The scale $s$ has factored out entirely. It factors out
because every term of the relation passes through $D$ exactly once, which makes the relation
homogeneous of degree one in $D$, and a homogeneous relation is invariant under $D \mapsto sD$ by
definition. Adding more relations of the same shape cannot help, since a system of relations each
invariant under that map has a scaling-invariant solution set. This is the part that is easy to get
wrong, and we got it wrong: it is tempting to conclude that no relation among the device's own
outputs can expose a scale error, and that an externally computed reference is therefore necessary.
It is not.

\section{Two ways to break the invariance}

A relation is blind to $s$ exactly when it is invariant under $D \mapsto sD$. Anything that breaks
that invariance exposes the scale error, and there are two ways to break it.

\paragraph{A term the device did not produce.} Put an independently computed value $V$ where a
device output stood. With one such anchor the residual on a probe $x$ is $(s-1)D(x) + \delta$. Two
qualifications matter. For a single probe this vanishes along a whole line in $(s, \delta)$, so one
probe separates nothing; it vanishes on every probe only if $s = 1$ and $\delta = 0$, and only if
the probes realise at least two distinct values of $D(x)$. And an anchor does not expose both errors
in general: for $\sum_i a_i D(x_i) + a_V V = 0$ the residual is $(1-s) a_V V + \delta \sum_i a_i$,
so a relation whose device terms are balanced sees $s$ and not $\delta$.

This route needs a reference, which is the thing relational testing exists to avoid.

\paragraph{A relation that mixes degrees.} The second route needs nothing external. The
multiplicative homomorphism
\[
D(a \boxtimes b) = D(a)\,D(b)
\]
is degree one on the left and degree two on the right. It consumes only device outputs, needs no
reference implementation and no stored answers, and under $D' = sD$ its residual is
\[
s\,D(ab) - s^2 D(a)D(b) \;=\; s(1-s)\,D(ab),
\]
which vanishes only at $s \in \{0, 1\}$. A relational suite can therefore catch a scale error. The
exact criterion for invariance is that every homogeneous component of the relation vanish
identically on correct decodes, of which same-degree homogeneity is the common sufficient case.

\section{A suite that missed it, and by how much}

We applied the criterion to a CKKS conformance suite of our own, released under
Apache-2.0~\cite{artifact}. It tests three laws: additive homomorphism, distributivity of plaintext
multiplication, and distributivity of ciphertext multiplication. Every one of them is degree one in
$D$, and two of the three are also balanced. By the criterion the suite is blind to every scale
error on all three laws, and blind to every offset on two of them. Measured on the software
reference at $N = 256$ and scale $2^{22}$, over five seeds:

\begin{center}
\begin{tabular}{@{}lccc@{}}
\toprule
decoder & add-hom & plainmul-dist & ctmul-dist \\
\midrule
honest                & 50.74 & 14.46 & 12.02 \\
scaled by $1.01$      & \textbf{50.76} & \textbf{14.46} & \textbf{12.02} \\
offset by $0.5$       & 2.36  & \textbf{14.74} & \textbf{12.41} \\
\bottomrule
\end{tabular}
\end{center}

The scale error is invisible on all three laws, and scores marginally \emph{higher} than honest on
the first. The offset is caught by additive homomorphism alone, which is the one law whose
coefficients do not sum to zero, and leaves the balanced pair unmoved. Both predictions hold
exactly. So the suite sees one of the two affine error modes, on one of its three laws, by accident
of which identities were chosen and not by design.

The same suite has a second mode in which the authority derives the probe plaintexts itself and
scores against them. That mode is the anchored case of the previous section, and it catches the
scale error on all three laws, at $6.64$ bits. The blindness is a property of testing relations
among the device's own outputs, not of the suite as a whole.

The law that would have caught the scale error is the multiplicative homomorphism, and it is one we
did not include. Implementing it on the same backend gives an honest score of $14.01$ bits and, at
$s = 1.01$, a score of $6.66$ bits, which matches the closed form $-\log_2(|s(1-s)|/s^2) = 6.66$ to
two decimals.

\section{Sensitivity is not refusal}

The last number is where the useful caution lies. A conformance suite does not report a score, it
issues a verdict, and a verdict compares the score against a floor. Under $D' = sD$ the achieved
precision of the multiplicative homomorphism is approximately $-\log_2(|1-s|/s)$, so the law refuses
a scale error only when
\[
|1 - s| \;\gtrsim\; s \cdot 2^{-\mathrm{floor}}.
\]
At the parameters of the example above the derived floor is $6.0$ bits, and $6.66 > 6.0$, so the law
detects the fault and still returns a pass. At $N = 8192$ and scale $2^{40}$, where the derived floor
is $19.0$ bits, the same $1\%$ error is refused comfortably.

Adding the right relation is therefore necessary and not sufficient. A suite reports what its
relations can see, and refuses what its floors put out of reach, and the two are set independently.

\section{What this means for relational FHE testing}

The criterion is not specific to our suite. It applies to any test that checks identities among a
device's own decoded outputs, which includes the metamorphic and equivalence-based designs cited
above. It does not apply to relation checks on ciphertext state, such as the key-free datapath
checks of~\cite{keyfree}, which verify fixed public linear maps by random projection and never
decode at all. The decoder is the layer this criterion is about. The practical form is short. For each relation in a suite, compute the coefficient sum and
the degree structure in $D$. A suite whose relations are all degree-one homogeneous is blind to
every decoder scale error, whatever else it tests and however many relations it has. A suite whose
relations are all balanced is blind to every constant decoder offset. Neither property is visible in
the identities as they are usually written, which is why a suite can have both and nobody notice.

We record this because we had both and did not notice.

\section*{Availability}

The suite, the faulty backends and the measurements above are at
\begin{center}
\url{https://github.com/VaultBytes/oracle-free-fhe-conformance}
\end{center}
under Apache-2.0.

\begin{thebibliography}{9}
\bibitem{eidolon} Z.~Xian, Z.~Yan, Y.~Chen, X.~Cao, F.~Ma, D.~Shi, Y.~Jiang. Eidolon: Perform
  Noise-Aware Fuzzing on FHE Libraries via Equivalence Expression Transformation. Proc. ACM Softw.
  Eng. 3(FSE), Article FSE102, 2026.
\bibitem{herta} Y.~Peng, D.~Xiao, Z.~Liu, Z.~Ji, S.~Wang. Detecting and Understanding
  Vulnerabilities in Fully Homomorphic Encryption Frameworks. arXiv:2606.22519, 2026.
\bibitem{ckks} J.~H. Cheon, A.~Kim, M.~Kim, Y.~Song. Homomorphic Encryption for Arithmetic of
  Approximate Numbers. ASIACRYPT 2017.
\bibitem{keyfree} B.~Alissaei. Where the Information Lives: A Key-Free Observability Decomposition
  for CKKS Hardware Monitors. Preprint, 20 June 2026. \url{https://doi.org/10.5281/zenodo.20774112}
\bibitem{artifact} Oracle-Free FHE Conformance.
  \url{https://github.com/VaultBytes/oracle-free-fhe-conformance}.
\end{thebibliography}

\end{document}
